\section{Metodología de implementación}

\subsection{Enfoque metodológico del proyecto}

La implementación del sistema de gestión en la nube para Machu Picchu Exclusive Tours E.I.R.L. se rige bajo un enfoque de transformación digital ágil, orientado a centralizar la data operativa y profesionalizar la atención al cliente internacional. Al tratarse de una solución basada en las funcionalidades estándar de Odoo Online (SaaS), el proyecto no requiere desarrollos de código a medida, lo que asegura una adopción inmediata, segura y de bajo riesgo técnico.

El enfoque principal es la transición del modelo manual (basado en Excel y WhatsApp) hacia un modelo digital inteligente, priorizando la seguridad de los documentos de identidad de los pasajeros y la trazabilidad de los itinerarios de viaje.

\subsection{Metodología ágil aplicada (Scrum)}

Se adopta la metodología Scrum adaptada al entorno operativo del sector turístico, lo que permite una ejecución por fases con validaciones continuas y mejora progresiva. Esta metodología facilita que el sistema se adapte de forma precisa a los flujos de la empresa mediante ciclos cortos de configuración y retroalimentación.

Las principales características de Scrum aplicadas al proyecto Machu Picchu Exclusive Tours son:

\begin{itemize}
    \item[A.] \textbf{Entregas parciales por módulos}: se presentó primero la configuración de CRM y leads.
    \item[b.] \textbf{Retroalimentación continua (feedback)}: ajuste inmediato de la interfaz tras la validación del cliente.
    \item[c.] \textbf{Flexibilidad en el alcance}: incorporación de nuevos requisitos funcionales externos al alcance.
\end{itemize}

\subsection{Herramientas de gestión y comunicación}

Para asegurar el cumplimiento de los hitos y mantener una documentación transparente, se utilizan las siguientes herramientas:

\begin{itemize}
    \item[a.] \textbf{Trello / Asana}: para el control visual de las tareas y el seguimiento semanal de los sprints.
    \item[b.] \textbf{Odoo Online}: instancia principal donde se realiza la parametrización de los módulos de CRM y Proyectos.
    \item[c.] \textbf{Canales de comunicación}: WhatsApp grupal del proyecto para coordinaciones inmediatas y reuniones de control semanales vía videollamada.
\end{itemize}

\subsection{Sprints de configuración}

A diferencia de un desarrollo tradicional, la metodología se centra en la configuración y adaptación de flujos. El trabajo se organiza en entregas parciales que permiten a la gerencia validar la usabilidad del CRM y el tablero de proyectos antes del despliegue final, asegurando que la herramienta responda a la realidad de tours que varían entre 2 a 20 días.

